\section{Related work and scope}\label{sec:related}

\subsection*{An apparatus that can question its own result}

Rheinberger's account of experimental systems helps frame the case
\cite{rheinberger1997epistemic}: apparatus makes questions answerable, but also
limits what can appear as an answer. Here the limit is concrete. A component
check could pass while its output had no live consumer. The response was to
inspect the handoff and construct a rejecting check, rather than strengthen
the prose around the passing result.

The pattern tradition records a recurring problem, its forces and a response
\cite{alexander1977a,gamma1994a,gabriel1writer}. Our PLoP companion presents the
mechanism catalogue; this paper follows the audit it prompted. The apparatus
patterns add a related question: what ensures that the mechanism described is
loaded, observed and independently assessed? Their APM field-trial record is
useful precisely because it includes misleading applications and subsequent
corrections. A promoted entry is not a universal effectiveness claim.

Winner's question about the politics of artifacts
\cite{winner1980artifacts} bears on the location of authority in this design.
The system may propose work and interpret observations, but adoption of a
preference or acceptance of an outcome cannot be silently inferred from the
same machinery that produced the proposal.

\subsection*{Agent work and operator participation}

Reflection and software-agent systems provide settings for execution and
feedback \cite{shinn2023reflexion,yang2024sweagent,wang2025openhands}.
The present contribution is narrower than a claim to a new general agent
architecture: it reports failures between declared mechanisms, recorded
execution and continuing operation in one stack. Neither an implementation
inventory nor a comparison of architectural diagrams establishes relative
performance.

Mixed-initiative design makes the cost of interruption and the allocation of
initiative explicit \cite{horvitz1999mixed}. Ostrom's action-arena perspective
similarly directs attention to participants, available actions and rules of
interaction \cite{ostrom2005understanding}. These perspectives help explain
why no-response cannot be treated as a transparent preference signal.
The paper identifies that confound and partial observations relevant to it;
it does not report a successful intervention on operator behaviour.

\subsection*{Coordination beyond the project}

Coase's account of coordination within firms, Beer's organisational cybernetics
and work on commons and peer production provide different settings for the
reporting question \cite{coase1937nature,beer1972brain,ostrom1990governing,
benkler2006wealth}. We use them to ask who supplies observations, who may demand
a report and who can accept an outcome. They do not establish that a particular
software instrument will be adopted or improve an organisation.

The empirical contribution remains the audit and its bounded demonstrations.
The issue board makes unresolved obligations visible; the operator model,
problem-level preferences and transfer study identify further tests. A useful
continuation will be a witnessed run that answers those tests, not a more
complete-looking description of the architecture.
